\title{Parameter-Efficient Orthogonal Finetuning via Factorization}

\begin{abstract}
The prevalence of large foundation models is rapidly increasing, yet the computational demands for training these models from the ground up are prohibitively high. Consequently, it is vital to develop methods for effectively customizing these high-capacity models for specific tasks. In this study, we focus on Orthogonal Finetuning (OFT), a structured and theoretically motivated approach for adapting models to downstream tasks. While OFT offers promising generalization, it is still burdened by the substantial number of trainable parameters necessitated by the dimensionality of orthogonal matrices. To mitigate this, we analyze OFT through the lens of information propagation, and articulate several essential criteria for achieving improved parameter-efficiency. Drawing inspiration from the Cooley-Tukey fast Fourier transform, which facilitates streamlined information processing, we introduce a compact orthogonal parameterization leveraging butterfly matrices. This parameterization is integrated with OFT to devise Orthogonal Butterfly~(BOFT), a new, highly parameter-efficient finetuning strategy. BOFT generalizes OFT by treating the latter as a specific instance, thus presenting a comprehensive framework for orthogonal finetuning. We systematically evaluate BOFT by adapting large-scale models, including vision transformers, language models, and text-to-image diffusion systems, across diverse tasks in vision and language domains.
\end{abstract}

\section{Introduction}

Models such as ChatGPT~\cite{brown2020language,chen2021evaluating} and Stable Diffusion~\cite{rombach2022high} have showcased the exceptional generalization capabilities characteristic of modern large foundation models. The extraordinary improvements in these models’ performance have coincided with significant growth in their parameter counts

(\emph{e.g.}, GPT-3 contains about 175B parameters). Consequently, training these models from the beginning is becoming an increasingly formidable challenge. Advancements in the broader research community therefore depend on methods that enable adapting existing foundation models to new tasks without full retraining. Efficient adaptation of pretrained models to downstream tasks is, therefore, imperative. Predominantly, three families of methods have emerged for this purpose: \emph{model finetuning}~\cite{hulora2022,zaken2022bitfit,guo2020parameter,raffel2020exploring,qiu2023controlling,zhang2022adaptive,chavan2023one}, which updates only a selected portion of a model’s parameters while preserving its architecture; \emph{adapter tuning}~\cite{rebuffi2017learning,houlsby2019parameter,pfeiffer2020adapterfusion,lin2020exploring,he2021towards}, in which supplemental trainable modules are introduced to the base network; and \emph{prompt tuning}~\cite{li2021prefix,lester2021power}, where trainable tokens are prepended to the input sequences. Of these, model finetuning stands out due to its conceptual simplicity, effectiveness, and, notably, because it avoids introducing any additional inference latency.

Model finetuning fundamentally aims to maintain similarity between the pretrained and finetuned models according to specific criteria, thereby retaining the knowledge acquired during pretraining. A commonly used tactic is to employ a low learning rate during finetuning, which heuristically limits the divergence between the original and updated model parameters. Considering the inherent structure of weight matrices, a more systematic strategy involves preserving the relationships within these matrices, particularly the pairwise angles among neurons. This concept underpins the development of Orthogonal Finetuning~(OFT)~\cite{qiu2023controlling}, a new finetuning methodology in which an identical orthogonal transformation is applied to all neurons within a given layer, ensuring the preservation of pairwise angular relationships throughout the finetuning procedure. While OFT has shown notable benefits in both generalizability and convergence, especially for finetuning text-to-image diffusion models~\cite{qiu2023controlling}, the approach can entail a significant number of trainable parameters owing to the large size of the orthogonal matrices involved. To mitigate this, OFT utilizes a block-diagonal orthogonal structure, reducing parameter count. Nonetheless, this increase in parameter efficiency introduces certain limitations—the orthogonal matrix must adhere to a rigid sparsity pattern, and orthogonal transformations are applied independently within each block. As a result, the block-diagonal design, while more efficient, may lead to less desirable inductive biases (\emph{e.g.}, diminished expressiveness and inability to approximate standard linear transformations), and critically, there remains an open question regarding optimal selection of the sparsity pattern.

A central challenge in this context is constructing a dense orthogonal matrix in a manner that remains efficient in terms of parameter count. Typically, forming a dense orthogonal matrix in $d$ dimensions entails managing $\mathcal{O}(d^2)$ parameters, making direct approaches seemingly impractical. We circumvent this constraint by introducing an alternative method: expressing the dense orthogonal matrix as a composition of several factorized sparse matrices. The motivation for this strategy arises from recognizing the possibility of decreasing the number of trainable parameters by opting for increased computational effort in lieu of additional storage. Specifically, as our representation relies on a sequence of sparse matrix multiplications to realize the orthogonal matrix, the requisite products must be recalculated at every training step. To frame the matrix factorization problem, we adopt an information transmission perspective—considering the synthesis of a dense orthogonal matrix as analogous to propagating information through a grid-based graph. Within this framework, the multiplication of sparse matrices to generate a dense orthogonal matrix corresponds to information flow across the graph, providing us with various opportunities to devise sparse factorizations that both constrain trainable parameters and retain expressive capacity sufficient to represent dense matrices.

To further enhance parameter efficiency within this transmission-based approach, we take inspiration from butterfly structures as used in the Cooley-Tukey fast Fourier transform algorithm~\cite{cooley1965algorithm}, where data can be rapidly communicated between different nodes~\cite{klasing1994broadcasting}. The butterfly graph intrinsic to the Cooley-Tukey algorithm gives rise to what is known as butterfly factorization, a specialized method for decomposing matrices sparsely. Using butterfly factorization, it becomes feasible to assemble a $d$-dimensional dense matrix from a product of $\mathcal{O}(\log d)$ sparse matrices, with each containing $\mathcal{O}(d)$ non-zero elements. If each such sparse matrix is also constructed to be orthogonal, the overall result is a highly efficient orthogonal parameterization requiring just $\mathcal{O}(d\log d)$ parameters—representing a dramatic reduction in complexity compared to conventional schemes. Leveraging this structure, we introduce Orthogonal Butterfly~(BOFT), a new approach for parameter-efficient finetuning. BOFT generalizes OFT, accommodating it as a specific instance, and lays out a comprehensive orthogonal finetuning paradigm. Notably, both the block-diagonal and butterfly constructions share a fundamental property: they enforce sparsity within the orthogonal matrices, thereby diminishing the number of parameters that must be learned. This makes for an illuminating comparison to methods such as LoRA~\cite{hulora2022}, which leverage low-rank matrices. Indeed, both low-rank and sparse matrices represent important subclasses of structured matrices~\cite{candes2011robust}, each promoting parameter efficiency through distinct structural constraints.

In contrast to the block-diagonal structure adopted by OFT for balancing model expressiveness and regularization, BOFT leverages the butterfly structure to establish a smoother continuum between matrices drawn from the full orthogonal group (high expressiveness) and identity matrices (strong regularization). This design permits the identification of a more compact hypothesis class situated within the orthogonal group, potentially improving downstream performance. Owing to the butterfly structure's prevalent role in efficient linear transformations—such as those underlying the discrete Fourier transform and the discrete cosine transform—we posit that this structural bias induces a beneficial inductive effect, fostering improved generalization while reducing the risk of overfitting. Our perspective stems from the observation that the butterfly structure can faithfully represent a range of canonical linear transforms, a feat that is unattainable for the block-diagonal construction utilized in OFT. Our key contributions can be summarized as follows:

\begin{itemize}[leftmargin=*,nosep]
    \item We address the challenge of parameter efficiency in orthogonal finetuning by developing a new information transmission framework. This framework reinterprets the construction of parameter-efficient dense orthogonal matrices as an information flow challenge across a grid-structured graph.
    \item Building on the butterfly patterns central to the Cooley-Tukey algorithm, we introduce Orthogonal Butterfly—a parameter-efficient orthogonal finetuning technique—within the scope of our proposed framework.
    \item We deliver theoretical analyses to elucidate how BOFT achieves a marked reduction in trainable parameter count while preserving robust expressiveness and stable training dynamics. Owing to its matrix factorization scheme, BOFT additionally supports a novel form of weight interpolation (see Figure~\ref{fig:boft_interpolation}).
    \item We are the first to extend orthogonal finetuning to a spectrum of tasks outside of controllable text-to-image generation~\cite{qiu2023controlling}, thereby revealing its versatility as a general model adaptation tool. Specifically, we apply BOFT to diverse adaptation scenarios in both computer vision and natural language processing, where it surpasses existing state-of-the-art methods by a wide margin, thus demonstrating its advanced parameter efficiency and generalization capability.
\end{itemize}

\section{Related Work}

\textbf{Parameter-efficient finetuning (PEFT)}. As foundation models (\emph{e.g.}, \cite{brown2020language,radford2021learning,rombach2022high,kirillov2023segment}) continue to expand in size and capability, the challenge of adapting these models to downstream applications in a manner that conserves parameters has become a focal point in the literature~\cite{treviso2023efficient,ding2023parameter,lialin2023scaling}. A wide variety of PEFT strategies exist~\cite{houlsby2019parameter,aghajanyan2020intrinsic,hulora2022,edalati2022krona,wang2022adamix,gheini2021cross,zaken2022bitfit,guo2020parameter,sung2021training,ansell2022composable,lester2021power,li2021prefix,vu2022spot,he2021towards,mao2021unipelt,karimi2021compacter,liu2022few,sung2022lst,chen2022parameter,jia2022visual,chen2022adaptformer,zhang2022neural,jie2023fact,lian2022scaling,luo2023towards}, among which reparameterization techniques~\cite{aghajanyan2020intrinsic,hulora2022,edalati2022krona} are especially pertinent to our research. LoRA~\cite{hulora2022} enhances a pretrained weight matrix by appending the product of two low-rank matrices, demonstrating strong results for a range of language processing benchmarks. Because LoRA constrains all supplementary matrices to the same rank, follow-up work \cite{zhang2022adaptive,valipour2022dylora,zhang2023increlora} explores layer-wise adaptive rank allocation to better utilize parameter resources. The ease of implementing such low-rank weight modifications has led to their widespread use \cite{dettmers2023qlora,zi2023delta,chavan2023one}. Drawing inspiration from research on hyperspherical energy and its relation to generalization~\cite{liu2018learning,liu2021orthogonal}, \cite{qiu2023controlling} introduces orthogonal finetuning, an alternative paradigm which directly learns orthogonal transformations of neuron activations within a layer; OFT demonstrates improved generalization and more reliable training trajectories than LoRA, particularly for text-to-image diffusion model finetuning. Notably, OFT generally incurs a larger number of trainable parameters compared to LoRA, motivating efforts to improve its parameter efficiency. Additionally, the generalizability of OFT across a broader range of adaptation scenarios (outside the domain of text-to-image diffusion~\cite{qiu2023controlling}) remains to be established. BOFT addresses the parameter overhead in OFT by employing butterfly factorization, thereby enabling orthogonal finetuning to be deployed in general adaptation contexts with enhanced efficiency.

\textbf{Butterfly structures}. The radix-2 Cooley-Tukey method decomposes an $N$-point discrete Fourier transform into smaller $\frac{N}{2}$-point transforms recursively, naturally giving rise to what is termed a butterfly structure; mathematically, this structure can be represented by products of sparse matrices known as butterfly matrices. Prior work has adopted butterfly matrices to parameterize orthogonal matrices—helping, for example, to avoid pivot operations in Gaussian elimination and improve computational speed~\cite{parker1995random}, to facilitate more robust training of recurrent neural networks~\cite{jing2017tunable}, and for use in kernel approximations~\cite{munkhoeva2018quadrature}. Additional applications include efficient learning of linear transforms via butterfly parameterizations~\cite{dao2019learning,dao2019kaleidoscope}, and accelerating neural network training through butterfly-structured layers~\cite{chen2022pixelated,dao2022monarch}. Moreover, butterfly structures benefit tasks such as rapid matrix-vector multiplication~\cite{michielssen1996multilevel,o2010algorithm}, compact matrix approximation~\cite{li2015butterfly}, and improved data transmission in networks~\cite{klasing1994broadcasting,li2003linear,rajasingh2016transmission}. Distinct from these prior applications, our work leverages butterfly structures specifically to advance the parameter efficiency within OFT.

\section{An Information Transmission View on Orthogonal Finetuning}

We begin by outlining key concepts underlying OFT. In order to adapt a pretrained weight matrix, OFT reformulates the updated weight matrix as a product between a learnable orthogonal matrix and the original, unmodified pretrained weights. Unlike LoRA—which introduces updates through an added low-rank matrix—OFT employs a multiplicative strategy, applying an orthogonal matrix to modify the existing weights. While LoRA attains parameter-efficiency using a low-rank formulation, the initial OFT approach instead leverages a block-diagonal orthogonal matrix~\cite{qiu2023controlling}; parameter efficiency increases as the diagonal blocks decrease in size. Figure~\ref{fig:comp_lora} offers a straightforward comparison of these methods. The primary rationale for using orthogonal transformations during weight matrix adaptation is to conserve the angular relationships between neuron vectors~\cite{liu2018learning,liu2021orthogonal,qiu2023controlling}, helping to retain much of the semantic information from pretraining.

More formally, for a pretrained linear layer $\bm{W}^0\in\mathbb{R}^{d\times n}$, OFT introduces a learnable orthogonal matrix $\bm{R}\in\mathbb{R}^{d\times d}$ and adjusts the computation in the forward pass from $\bm{z} = (\bm{W}^0)^\top \bm{x}$ to $\bm{z} = (\bm{R} \bm{W}^0)^\top \bm{x}$, with $\bm{x}\in\mathbb{R}^d$ representing the input and $\bm{z}\in\mathbb{R}^n$ the output vector. To begin with a zero-update, $\bm{R}$ is initialized as the identity matrix. During fine-tuning, to maintain the orthogonality of $\bm{R}$, we adopt the Cayley parameterization as in \cite{qiu2023controlling,liu2021orthogonal}, i.e., $\bm{R}=(\bm{I}+\bm{Q})(\bm{I}-\bm{Q})^{-1}$, where $\bm{Q}$ is a skew-symmetric matrix so that $\bm{Q}=-\bm{Q}^\top$. To further economize on parameters, the block-diagonal design is employed by defining $\bm{R}$ as $\text{diag}(\bm{R}_1,\bm{R}_2,\cdots,\bm{R}_r)$, where each $\bm{R}_i\in\mathcal{R}^{b\times b}$ (for $i=1,...,r$) is a smaller orthogonal matrix and $br=d$. However, this parameter-saving tactic introduces an implicit assumption: the dimensions of each neuron (corresponding to columns in $\bm{W}^0$) are partitioned into $r$ distinct groups, with each group being independently transformed by its associated orthogonal matrix. While block-diagonal orthogonal structures have demonstrated empirical success, the practice of partitioning neuron dimensions strictly by index lacks theoretical grounding, underscoring the potential value of dense orthogonal matrices. This prompts a fundamental inquiry: \emph{Is it possible to devise a dense orthogonal matrix that maintains parameter efficiency?}

To explore this problem, we suggest constructing a dense orthogonal matrix by multiplying several sparse orthogonal matrices. To systematically investigate the sparsity configurations in orthogonal matrix decompositions, we reinterpret the task of generating a dense orthogonal matrix in OFT as analogous to an information transmission scenario. Concretely, the creation of a dense matrix $\bm{R}\in\mathbb{R}^{d\times d}$ via the product $\thickmuskip=2mu \medmuskip=2mu \bm{R}=\bm{B}_m\bm{B}_{m-1}\cdots\bm{B}_1$, with $m$ square matrices, can be recast as a process of communicating information across a $d\times (m+1)$ grid of nodes (see Figure~\ref{fig:info_trans} for reference). This perspective arises from noticing that a dense square matrix of size $d$ can represent full connectivity between two sets of $d$ nodes. In this context, the value of $R_{ij}$ in $\bm{R}$ determines whether information is transmitted from node $j$ to node $i$: a zero entry implies a blockade, while a non-zero entry allows passage of information. Thus, the factorization $\bm{R} = \bm{B}_m \bm{B}_{m-1}\cdots \bm{B}_1$ can be interpreted as sequential exchanges of information as dictated by the connectivity graphs defined by each $\bm{B}_i$. Information traverses according to the connections in $\bm{B}_1$ at the outset and propagates through to $\bm{B}_n$ at the end. For instance, Figure~\ref{fig:info_trans} illustrates the case where $\bm{R} = \bm{B}_5\bm{B}_4\bm{B}_3\bm{B}_2\bm{B}_1$, displaying both the sparsity patterns and their resulting graph structure. This graph reflects the sequential unfolding of the matrix product. Within the graph, $\bm{B}_i$ acts as the connection matrix linking nodes from the $i$-th to the $(i+1)$-th level; specifically, the $(j_1, j_2)$ element of $\bm{B}_i$ encodes a directed edge from node $j_2$ at level $i$ to node $j_1$ at level $(i+1)$, with zero entries indicating absence of an edge. To ensure that the product $\bm{B}_5\bm{B}_4\bm{B}_3\bm{B}_2\bm{B}_1$ produces a dense $\bm{R}$, it is required that information from every node at the first level can reach every node at the sixth level. When restricting to $\bm{R} = \bm{B}_4\bm{B}_3\bm{B}_2\bm{B}_1$ (involving levels one and five), we observe that information originating from node $1$ fails to reach node $3$, corresponding to a zero at position $(3,1)$ in the sparsity pattern of $\bm{B}_4\bm{B}_3\bm{B}_2\bm{B}_1$. The analogy between the graph representation and matrix sparsity persists in the cases of shorter products such as $\bm{R} = \bm{B}_3\bm{B}_2\bm{B}_1$ and $\bm{R} = \bm{B}_2\bm{B}_1$. In summary, achieving a dense matrix $\bm{R}\in\mathbb{R}^{d\times d}$ requires that the sequence of factors $\bm{B}_m, \cdots, \bm{B}_1$ specifies a system of directed edges on a $d\times (m+1)$ grid, with each edge only bridging adjacent levels (columns), ensuring the possibility of information transfer from any starting node in the first level to every endpoint in the $(m+1)$-th level.

The perspective of information transmission offers valuable insights into the sparsity characteristics of factorization matrices within OFT. As shown in Figure~\ref{fig:blk_diag}, $\bm{R}$ in the original OFT displays a block-diagonal structure. While this configuration effectively reduces the total number of parameters requiring optimization, it fails to represent a fully dense matrix $\bm{R}$. Our objective is to construct a dense orthogonal matrix by combining $m$ sparse orthogonal matrices, aiming to minimize the count of non-redundant trainable parameters. Adopting the information transmission approach, we articulate two primary criteria: \emph{(i)} \textbf{dense connectivity}, ensuring that every node at the initial layer is linked, via at least one path, to all nodes at the terminal layer; and \emph{(ii)} \textbf{minimum free edges}, which seeks the lowest possible number of edges consistent with the orthogonality requirement. Orthogonality imposes a nuanced restriction on the allowable edges between successive layers. Specifically, maintaining full rank for each matrix $\bm{B}_i$ (an essential property for orthogonality) necessitates $d$ edges to establish a bijective mapping between the $i$-th and $(i=1)$-th layers. Consequently, there must be at least $d$ edges joining consecutive layers (such as 4 edges in the illustration in Figure~\ref{fig:info_trans}). Since these $d$ essential connections are required for orthogonality and do not involve trainable parameters (for instance, a $d\times d$ orthogonal matrix with exactly $d$ non-zeros restricts them to values of $\pm 1$), they are not included when counting the number of edges. Compared to arbitrary full matrices, orthogonal matrices necessitate fewer trainable parameters, causing additional dependencies among the present edges. As an illustration, in a $2\times 2$ orthogonal matrix, a single zero in position $(1,1)$ implies another zero in position $(2,2)$—indicating that eliminating one edge may compel the elimination of another. Therefore, the orthogonality property may force modifications in the set of edges for each valid connection configuration. Visualizing the nonzero pattern in the resulting orthogonal matrix underscores the utility of the information transmission perspective in the context of OFT, as our primary interest lies in the distribution and positions of trainable, nonzero entries of $\bm{R}$, independent of their actual numerical values.

A straightforward approach to creating dense connections between two layers requires $\mathcal{O}(d^2)$ edges, corresponding to one dense orthogonal matrix and yielding $d^2-d$ edges (for instance, with $d=4$, there are 12 connections). Figure~\ref{fig:info_trans} illustrates a specific instance of matrix factorization that only utilizes $10$ edges, which is fewer than what is needed by a full dense orthogonal matrix. This conceptual model allows for the examination of the parameter-efficiency of OFT from a graph-theoretic viewpoint, facilitating the easy construction of efficient factorizations. Motivated by a distinctive topology used in the Cooley-Tukey algorithm, known as butterfly graphs~\cite{cooley1965algorithm}, we observe that these structures can effectively connect $d$ source nodes to $d$ destination nodes using just $\mathcal{O}(d\log d)$ edges. As an example, the connectivity arrangement in Figure~\ref{fig:info_trans} requires $10$ edges, whereas the butterfly topology achieves full connectivity with only $8$ edges. In the following, we describe how leveraging the butterfly structure can enhance parameter efficiency.

\section{Orthogonal Parameterization by Butterfly Factorization}

The butterfly structure, first introduced in the Cooley-Tukey algorithm for accelerating the fast Fourier transform, is relevant to our context. In the Fourier transform, a small modification in the frequency domain affects the entire spatial domain, which resembles the situation in information transmission where each node at the input layer can influence every output node. Moreover, the butterfly structure has been widely adopted in computer network designs~\cite{solihin2015fundamentals,li2011linear} due to its efficacy in facilitating information exchange. Assuming $k\geq 2$ is a power of $2$, we define the \emph{butterfly factor} $\bm{B}^F(k)$ as follows:

\begin{equation}\label{eq:but_factor}
\begin{aligned}
    \bm{B}^F(k)=\begin{bmatrix}
    \text{diag}(\bm{d}_1) & \text{diag}(\bm{d}_2) \\
    \text{diag}(\bm{d}_3) & \text{diag}(\bm{d}_4)
    \end{bmatrix}\in\mathbb{R}^{k\times k},
\end{aligned}
\end{equation}

where for each $i$, $\bm{d}_i\in\mathbb{R}^{\frac{k}{2}}$ are vectors. For $d=2^N$, we recursively construct the $d$-dimensional \emph{butterfly matrix} $\bm{B}(d)\in\mathbb{R}^{d\times d}$ by

\begin{equation}
\begin{aligned}
    \!\!\bm{B}(d)=\tilde{\bm{B}}(d,d)\!\cdot\!\begin{bmatrix}
    \bm{B}_1(\frac{d}{2})\!\!\!\!\! & \bm{0} \\
    \bm{0}\!\!\!\!\! &\bm{B}_2(\frac{d}{2})
    \end{bmatrix}= \tilde{\bm{B}}(d,d)\tilde{\bm{B}}(d,\frac{d}{2})\cdots\tilde{\bm{B}}(d,2),
\end{aligned}
\end{equation}

Here, $\bm{B}_1(\frac{d}{2})$ and $\bm{B}_2(\frac{d}{2})$ denote two butterfly matrices, each of dimension $\frac{d}{2}$. We introduce the \emph{butterfly component} as $\thickmuskip=2mu \medmuskip=2mu \tilde{\bm{B}}(d,k)=\text{diag}(\bm{B}^F_1(k),\cdots,\bm{B}^F_{d/k}(k))$, which forms a $d\times d$ block-diagonal matrix with individual blocks of size $k$. The matrices $\bm{B}^F_i(k),\forall i$ represent the butterfly factors as specified in Equation~\ref{eq:but_factor}. This framework allows us to use the butterfly matrix for orthogonal matrix parameterization, provided that every multiplicative factor $\tilde{\bm{B}}(d,k)$ within $\bm{B}(d)$ maintains orthogonality. Focusing first on the block-diagonal case $\tilde{\bm{B}}(d,2)$, where each block has size $2$, orthogonality can be imposed by using the Cayley transform or a $2$-dimensional rotation for parameterizing each block, as outlined in \cite{liu2021orthogonal,qiu2023controlling}. Each butterfly component exhibits a non-zero structure that can be considered as a permuted version of the pattern in $\tilde{\bm{B}}(d,2)$, facilitating straightforward orthogonal parameterizations for all components. This approach leads to an efficient scheme for parameterizing orthogonal matrices using multiple $2\times 2$ orthogonal matrices as building blocks. Drawing from \cite{chen2022pixelated}, we extend butterfly matrices and introduce a block butterfly component $\tilde{\bm{B}}^b(d,k)$, wherein every entry in each vector $\bm{d}_i,\forall i$ is a $b\times b$ matrix. Ensuring the orthogonality of the block butterfly component $\tilde{\bm{B}}^b(d,2)$ is achieved by parameterizing each constituent $2b\times 2b$ block as an orthogonal matrix. For higher block sizes ($k>2$), the non-zero layout in $\tilde{\bm{B}}^b(d,k)$ corresponds to a blockwise permutation of the pattern from $\tilde{\bm{B}}^b(d,2)$, which in turn allows similar orthogonal parameterizations. Bringing these components together, the BOFT forward pass can be described as follows:

\begin{equation}
    \begin{aligned}\nonumber
    \bm{z}=\big(\bm{R}(m,b)\cdot\bm{W}^0\big)^\top\bm{x},~~~~\text{s.t.}~\bigg{\{}\bm{R}(m,b)=\prod_{i=1}^m \tilde{\bm{B}}^b_{(d,i)}~~\&~~\underbrace{\big(\tilde{\bm{B}}^b_{(d,j)}\big)^\top\tilde{\bm{B}}^b_{(d,j)}=\tilde{\bm{B}}^b_{(d,j)}\big(\tilde{\bm{B}}^b_{(d,j)}\big)^\top\!=\bm{I}_d}_{\forall j\in [1, m]}\bigg{\}},
    \end{aligned}
\end{equation}

Here, the notation $\tilde{\bm{B}}^b(d,2^{m-i+1})$ has been simplified to $\tilde{\bm{B}}^b_{(d,i)}$, and $\bm{I}_d$ refers to the $d \times d$ identity matrix. The orthogonal transformation $\bm{R}(m,b)\in\mathbb{R}^{d\times d}$ is assembled as the product of several orthogonal butterfly matrices. We designate BOFT constructed with $\bm{R}(m, \frac{b}{2})$ as BOFT($m$,$b$), assuming $b\geq 2$. Notably, when $\thickmuskip=2mu \medmuskip=2mu m=1$, BOFT($1$,$b$) is equivalent to the block-diagonal OFT described in~\cite{qiu2023controlling} using blocks of size $b$. When $b=d$, BOFT($1$,$d$) corresponds to the standard OFT~\cite{qiu2023controlling}, using a fully unconstrained orthogonal matrix. Additionally, choosing $m=\log \frac{2d}{b}$ and $b$ yields BOFT($\log \frac{2d}{b}$,$b$), which employs the block butterfly matrix $\bm{B}^{\frac{b}{2}}(d)$ as $\bm{R}$ and results in a dense orthogonal transformation. More generally, BOFT($m$,$b$) incorporates $\frac{1}{2}(b-1)dm$ independent parameters when adapting a $d\times n$ linear layer. For the scenario with butterfly matrices, i.e., $m=\log d$ and $b=2$, BOFT operates with $\mathcal{O}(d\log d)$ parameters. By contrast, the original OFT with a dense orthogonal matrix requires $\mathcal{O}(d^2)$ parameters, while the block-diagonal variant with $r$ blocks uses $\mathcal{O}(bd)$. As a result, in the conventional OFT, one must use a block size $b=d$ to generate a dense orthogonal mapping, but BOFT achieves this flexibility for any $b$.

\textbf{BOFT with Identity Initialization}. For parameter-efficient tuning approaches, initializing from the precise weights of the pretrained model is crucial so that the adjusted model remains close to its original state. For instance, LoRA initializes low-rank weight updates as zeros. In the context of BOFT, every butterfly module is set as an identity at initialization (i.e., the Cayley-transformed skew-symmetric matrices start from all-zeros).

\textbf{Multiplicative Dropout in BOFT}. LoRA~\cite{hulora2022} employs a Dropout layer on the low-rank adaptation weights to curb overfitting. While traditional Dropout~\cite{srivastava2014dropout} is compatible with LoRA’s additive updates, it is unsuitable for BOFT due to the multiplicative fashion in which weights are updated. To address this limitation, we introduce a multiplicative Dropout mechanism specifically for BOFT. Since $\bm{R}(m,b)$ is constructed from $m$ butterfly matrices that can be restructured into $2b\times 2b$ block-diagonal orthogonal forms, the multiplicative Dropout works as follows: a proportion $p_1$ of butterfly matrices and a proportion $p_2$ of the diagonal blocks within each butterfly are randomly selected and then replaced with identity matrices.

\section{Intriguing Insights and Discussions}

\textbf{Expressivity of BOFT}. The butterfly architecture, together with appropriate permutations, is capable of exactly representing numerous well-known fast linear transforms~\cite{dao2019learning,dao2019kaleidoscope} (such as the fast Fourier transform and the Hadamard transform). Nonetheless, the degree to which our orthogonal butterfly matrix can approximate arbitrary orthogonal matrices remains an open question. To investigate this, we perform experiments that fit a randomly sampled dense orthogonal matrix~\cite{anderson1987generation} of dimension $512\times 512$. The findings, depicted in Figure~\ref{fig:para_eff}, are averaged over 10 distinct random initializations. Here, the vertical axis represents the error in the approximation, while the horizontal axis shows the count of effective trainable parameters. Each curve with matching color corresponds to a BOFT variant sharing the same block size. The leftmost point on each curve is the error of BOFT($1$,$b$), i.e., a purely block-diagonal OFT structure of block size $b$. BOFT tends to provide superior parameter efficiency over OFT; for instance, BOFT($9$,$2$) is more expressive than BOFT($1$,$16$) yet requires fewer parameters. In general, choosing a smaller $b$ alongside a larger $m$ in BOFT configurations leads to higher parameter efficiency. For example, BOFT($6$,$4$) achieves comparable approximation error to BOFT($2$,$16$) while utilizing significantly fewer parameters. The butterfly construction, by design, embodies a more structured subset within the orthogonal group relative to the block-diagonal approach, thereby making BOFT demonstrably more expressive than OFT with equivalent block size.

\begin{theorem}[Expressivity of BOFT]\label{thm:exp_boft}
    BOFT achieves higher expressivity compared to OFT for a given block size. It is possible to approximate any orthogonal matrix of size $d$ via a product of butterfly matrices: $\bm{B}_{d-1,1}(d)\bm{B}^\top_{d-1,2}(d)\cdots\bm{B}_{1,1}(d)\bm{B}^\top_{1,2}(d)$, with each $\bm{B}_{i,j}(d)$ representing a butterfly matrix.
\end{theorem}

The implication of Theorem~\ref{thm:exp_boft} is a straightforward generalization for BOFT: the resulting orthogonal matrix can be written as $\thickmuskip=2mu \medmuskip=2mu \bm{R}^G(m_1,b_1,m_2,b_2,l)=\bm{R}_{l,1}(m_1,b_1)\bm{R}_{l,2}^T(m_2,b_2)\cdots\bm{R}_{1,1}(m_1,b_1)\bm{R}_{1,2}^T(m_2,b_2)$, in which $\bm{R}_{i,j}^T(m,b)$ denotes an orthogonal matrix employed in BOFT. Specifically, when $m_1=m_2=\log d$, $b_1=b_2=2$, and $l=d-1$, the construction $\bm{R}^G(m_1,b_1,m_2,b_2,l)$ is capable of spanning the full orthogonal group. This recursive formulation is also referred to as the kaleidoscope hierarchy~\cite{dao2019kaleidoscope}. It should be emphasized, however, that greater expressiveness does not necessarily guarantee improved finetuning outcomes. For example, although full finetuning is maximally expressive in theory, its practical performance can be suboptimal. The fundamental challenge lies in balancing expressivity and regularization to promote model generalizability during finetuning. By imposing structured constraints, BOFT expands the parameter space for finetuning and enables improved trade-offs between model expressivity and regularization.

\textbf{Spectral properties}. Compared to LoRA, orthogonal finetuning more effectively maintains the spectral characteristics of the pretrained weight matrix $\bm{W}^0$, as it exactly conserves its spectral norm. This is evident through the singular value decomposition: $\bm{W}^0 = \bm{U}\bm{\Sigma}\bm{V}^\top$, where $\bm{U}$ and $\bm{V}$ are orthogonal matrices and $\bm{\Sigma}$ contains the singular values along its diagonal. In both OFT and BOFT, the finetuning process involves left-multiplication by an orthogonal matrix $\bm{R}$, resulting in new weights $\bm{R}\bm{U}\bm{\Sigma}\bm{V}^\top$. This operation leaves the largest singular value unchanged, preserving the spectral norm of $\bm{W}^0$. Maintaining the spectral norm in this manner has been demonstrated to enhance both training robustness and model generalization~\cite{miyato2018spectral, yoshida2017spectral}. Further mathematical details can be found in Appendix~\ref{app:sn_property}.

\textbf{Orthogonal finetuning as learning bilinear similarity}. The BOFT framework can be interpreted as optimizing a bilinear similarity function of the form $\bm{w}^0_i\bm{R}\bm{x}$, where $\bm{w}^0_i$ refers to the $i$-th column (i.e., neuron) of the weight matrix $\bm{W}^0$. This perspective casts BOFT as a learning procedure for the bilinear similarity matrix $\bm{R}$ under the orthogonality constraint, a connection that is closely related to concepts in distance metric learning~\cite{xing2002distance} and the theory of bilinear forms~\cite{roman2005advanced}.

\textbf{Inductive bias and generalization in BOFT}. The structure of $\bm{R}(m, b)$ in BOFT often corresponds to a specific subset of the orthogonal group, effectively reducing the size of the hypothesis space and thus introducing an inductive bias. The inductive bias brought about by butterfly factorization is argued to be favorable for generalization, as this pattern appears in a variety of fundamental linear transforms~\cite{dao2019kaleidoscope}, including the discrete Fourier, sine/cosine, and Hadamard transforms. Additionally, the inherent sparsity from matrix factorization in BOFT may introduce further implicit inductive biases~\cite{gunasekar2017implicit, li2018algorithmic, liu2019neural}.

\textbf{Relation to butterfly-based sparse training}. Several studies~\cite{dao2019learning,dao2019kaleidoscope,chen2022pixelated,dao2022monarch} have investigated sparse training methodologies leveraging butterfly parameterizations. These prior works primarily emphasize the direct reparameterization of weight matrices using the butterfly structure and train neural networks from their initial state. In \cite{dao2022monarch}, the approach entails projecting pretrained weights onto a variation of butterfly matrices before finetuning the resulting components for downstream applications. In contrast, BOFT introduces an alternative finetuning approach, where layer-shared weight matrices are utilized to transform the model's weights.

\input{rebuttal-results}

\section{Concluding Remarks and Limitations}

This study introduces Orthogonal Butterfly, a universal parameter-efficient finetuning approach for foundation models that capitalizes on the butterfly structure. The central concept underpinning its parameter-efficiency lies in constructing a dense orthogonal matrix as a product of multiple sparse orthogonal matrices. To facilitate the identification of suitable matrix factorizations, we develop a graphical information transmission paradigm. Within this paradigm, we determine that the butterfly structure offers an effective means for sparse orthogonal matrix factorization aligned with our objectives. Through extensive experiments, we provide evidence of BOFT’s strong empirical performance across finetuning tasks involving large language models, extensive vision models, and text-to-image generative systems. Our findings corroborate the broad applicability and superiority of BOFT as a general finetuning strategy.

Nonetheless, BOFT is not devoid of limitations. Given that the final orthogonal matrix in BOFT is constructed by multiplying several orthogonal matrices, the computational overhead during training is marginally higher than that of OFT. Enhancing the training efficiency of BOFT therefore represents a pertinent avenue for future work. Notably, once finetuning is complete, the orthogonal matrices produced by BOFT can be seamlessly integrated into the pretrained model, resulting in no increase in inference time. Additionally, it remains unclear if the butterfly topology is the optimal choice for information transmission. The proposed information transmission framework opens the door for insights from computer networking—a domain where network topologies and their information transmission efficiencies are extensively investigated. This suggests that even more effective network architectures could potentially be leveraged for constructing dense orthogonal matrices.